% Arquitectura del autoencoder y del embedding (espacio latente).
%
% Fragmento para % Arquitectura del autoencoder y del embedding (espacio latente).
%
% Fragmento para % Arquitectura del autoencoder y del embedding (espacio latente).
%
% Fragmento para % Arquitectura del autoencoder y del embedding (espacio latente).
%
% Fragmento para \input{Figures/autoencoder_embedding}.
% Requiere en el preambulo del documento principal:
%   \usepackage{tikz}
%   \usetikzlibrary{decorations.pathreplacing,arrows.meta}
%   \usepackage{amsmath,amssymb}
%
% x     : vector de entrada de dimension n.
% z     : embedding (representacion latente) de dimension d << n.
% xhat  : reconstruccion de la entrada.
% E, D  : encoder y decoder, cada uno una MLP que reduce/expande la dimension.

\begin{figure}[htbp]
  \centering
  \begin{tikzpicture}[
      >={Latex[length=2.2mm]},
      vector/.style   = {draw=black!75, fill=white, line width=0.6pt},
      celda/.style    = {draw=black!35, line width=0.3pt},
      bloque/.style   = {draw=black!75, line width=0.6pt},
      encoder/.style  = {bloque, fill=black!8},
      decoder/.style  = {bloque, fill=black!8},
      latente/.style  = {draw=black!75, fill=black!18, line width=0.6pt},
      texto/.style    = {font=\footnotesize, align=center},
      nombre/.style   = {font=\small, align=center},
      flecha/.style   = {->, draw=black!75, line width=0.7pt},
    ]
    \useasboundingbox (-1.7,-4.5) rectangle (12.2,4.6);

    %------------------------------------------------ vector de entrada
    \draw[vector] (0,-2.2) rectangle (0.5,2.2);
    \foreach \y in {-1.65,-1.1,...,1.75}
      \draw[celda] (0,\y) -- (0.5,\y);
    \node[nombre, anchor=south] at (0.25,2.6) {$\mathbf{x}\in\mathbb{R}^{n}$};
    \node[texto, anchor=north] at (0.25,-2.6) {entrada};

    %------------------------------------------------ encoder
    \draw[encoder] (1.4,-2.2) -- (4.3,-0.75) -- (4.3,0.75) -- (1.4,2.2) -- cycle;
    \node[nombre] at (2.85,0) {$E_{\phi}$};
    \node[texto, anchor=north] at (2.85,-2.6) {codificador};

    %------------------------------------------------ espacio latente
    \draw[latente] (5.2,-0.75) rectangle (5.7,0.75);
    \foreach \y in {-0.375,0,0.375}
      \draw[celda] (5.2,\y) -- (5.7,\y);
    \node[nombre, anchor=south] at (5.45,1.15) {$\mathbf{z}\in\mathbb{R}^{d}$};
    \node[texto, anchor=north] at (5.45,-1.15)
      {embedding\\[-1pt] $d\ll n$};

    %------------------------------------------------ decoder
    \draw[decoder] (6.6,-0.75) -- (9.5,-2.2) -- (9.5,2.2) -- (6.6,0.75) -- cycle;
    \node[nombre] at (8.05,0) {$D_{\theta}$};
    \node[texto, anchor=north] at (8.05,-2.6) {decodificador};

    %------------------------------------------------ vector de salida
    \draw[vector] (10.4,-2.2) rectangle (10.9,2.2);
    \foreach \y in {-1.65,-1.1,...,1.75}
      \draw[celda] (10.4,\y) -- (10.9,\y);
    \node[nombre, anchor=south] at (10.65,2.6) {$\hat{\mathbf{x}}\in\mathbb{R}^{n}$};
    \node[texto, anchor=north] at (10.65,-2.6) {reconstrucci\'on};

    %------------------------------------------------ flujo
    \draw[flecha] (0.6,0)  -- (1.3,0);
    \draw[flecha] (4.4,0)  -- (5.1,0);
    \draw[flecha] (5.8,0)  -- (6.5,0);
    \draw[flecha] (9.6,0)  -- (10.3,0);

    %------------------------------------------------ llaves de las etapas
    \draw[decorate, decoration={brace, amplitude=5pt}, draw=black!60]
      (1.4,2.75) -- (4.3,2.75);
    \node[texto, anchor=south] at (2.85,3.15)
      {compresi\'on\\[-1pt] $n\to d$};
    \draw[decorate, decoration={brace, amplitude=5pt}, draw=black!60]
      (6.6,2.75) -- (9.5,2.75);
    \node[texto, anchor=south] at (8.05,3.15)
      {expansi\'on\\[-1pt] $d\to n$};

    %------------------------------------------------ funcion de perdida
    \node[texto, anchor=north] at (5.45,-3.35)
      {$\mathcal{L}(\phi,\theta)=\lVert\mathbf{x}-\hat{\mathbf{x}}\rVert^{2}$};
  \end{tikzpicture}
  \caption{Esquema del autoencoder: el codificador comprime la entrada en un
    embedding de baja dimensión y el decodificador reconstruye la entrada a
    partir de él.}
  \label{fig:autoencoder-embedding}
\end{figure}
.
% Requiere en el preambulo del documento principal:
%   \usepackage{tikz}
%   \usetikzlibrary{decorations.pathreplacing,arrows.meta}
%   \usepackage{amsmath,amssymb}
%
% x     : vector de entrada de dimension n.
% z     : embedding (representacion latente) de dimension d << n.
% xhat  : reconstruccion de la entrada.
% E, D  : encoder y decoder, cada uno una MLP que reduce/expande la dimension.

\begin{figure}[htbp]
  \centering
  \begin{tikzpicture}[
      >={Latex[length=2.2mm]},
      vector/.style   = {draw=black!75, fill=white, line width=0.6pt},
      celda/.style    = {draw=black!35, line width=0.3pt},
      bloque/.style   = {draw=black!75, line width=0.6pt},
      encoder/.style  = {bloque, fill=black!8},
      decoder/.style  = {bloque, fill=black!8},
      latente/.style  = {draw=black!75, fill=black!18, line width=0.6pt},
      texto/.style    = {font=\footnotesize, align=center},
      nombre/.style   = {font=\small, align=center},
      flecha/.style   = {->, draw=black!75, line width=0.7pt},
    ]
    \useasboundingbox (-1.7,-4.5) rectangle (12.2,4.6);

    %------------------------------------------------ vector de entrada
    \draw[vector] (0,-2.2) rectangle (0.5,2.2);
    \foreach \y in {-1.65,-1.1,...,1.75}
      \draw[celda] (0,\y) -- (0.5,\y);
    \node[nombre, anchor=south] at (0.25,2.6) {$\mathbf{x}\in\mathbb{R}^{n}$};
    \node[texto, anchor=north] at (0.25,-2.6) {entrada};

    %------------------------------------------------ encoder
    \draw[encoder] (1.4,-2.2) -- (4.3,-0.75) -- (4.3,0.75) -- (1.4,2.2) -- cycle;
    \node[nombre] at (2.85,0) {$E_{\phi}$};
    \node[texto, anchor=north] at (2.85,-2.6) {codificador};

    %------------------------------------------------ espacio latente
    \draw[latente] (5.2,-0.75) rectangle (5.7,0.75);
    \foreach \y in {-0.375,0,0.375}
      \draw[celda] (5.2,\y) -- (5.7,\y);
    \node[nombre, anchor=south] at (5.45,1.15) {$\mathbf{z}\in\mathbb{R}^{d}$};
    \node[texto, anchor=north] at (5.45,-1.15)
      {embedding\\[-1pt] $d\ll n$};

    %------------------------------------------------ decoder
    \draw[decoder] (6.6,-0.75) -- (9.5,-2.2) -- (9.5,2.2) -- (6.6,0.75) -- cycle;
    \node[nombre] at (8.05,0) {$D_{\theta}$};
    \node[texto, anchor=north] at (8.05,-2.6) {decodificador};

    %------------------------------------------------ vector de salida
    \draw[vector] (10.4,-2.2) rectangle (10.9,2.2);
    \foreach \y in {-1.65,-1.1,...,1.75}
      \draw[celda] (10.4,\y) -- (10.9,\y);
    \node[nombre, anchor=south] at (10.65,2.6) {$\hat{\mathbf{x}}\in\mathbb{R}^{n}$};
    \node[texto, anchor=north] at (10.65,-2.6) {reconstrucci\'on};

    %------------------------------------------------ flujo
    \draw[flecha] (0.6,0)  -- (1.3,0);
    \draw[flecha] (4.4,0)  -- (5.1,0);
    \draw[flecha] (5.8,0)  -- (6.5,0);
    \draw[flecha] (9.6,0)  -- (10.3,0);

    %------------------------------------------------ llaves de las etapas
    \draw[decorate, decoration={brace, amplitude=5pt}, draw=black!60]
      (1.4,2.75) -- (4.3,2.75);
    \node[texto, anchor=south] at (2.85,3.15)
      {compresi\'on\\[-1pt] $n\to d$};
    \draw[decorate, decoration={brace, amplitude=5pt}, draw=black!60]
      (6.6,2.75) -- (9.5,2.75);
    \node[texto, anchor=south] at (8.05,3.15)
      {expansi\'on\\[-1pt] $d\to n$};

    %------------------------------------------------ funcion de perdida
    \node[texto, anchor=north] at (5.45,-3.35)
      {$\mathcal{L}(\phi,\theta)=\lVert\mathbf{x}-\hat{\mathbf{x}}\rVert^{2}$};
  \end{tikzpicture}
  \caption{Esquema del autoencoder: el codificador comprime la entrada en un
    embedding de baja dimensión y el decodificador reconstruye la entrada a
    partir de él.}
  \label{fig:autoencoder-embedding}
\end{figure}
.
% Requiere en el preambulo del documento principal:
%   \usepackage{tikz}
%   \usetikzlibrary{decorations.pathreplacing,arrows.meta}
%   \usepackage{amsmath,amssymb}
%
% x     : vector de entrada de dimension n.
% z     : embedding (representacion latente) de dimension d << n.
% xhat  : reconstruccion de la entrada.
% E, D  : encoder y decoder, cada uno una MLP que reduce/expande la dimension.

\begin{figure}[htbp]
  \centering
  \begin{tikzpicture}[
      >={Latex[length=2.2mm]},
      vector/.style   = {draw=black!75, fill=white, line width=0.6pt},
      celda/.style    = {draw=black!35, line width=0.3pt},
      bloque/.style   = {draw=black!75, line width=0.6pt},
      encoder/.style  = {bloque, fill=black!8},
      decoder/.style  = {bloque, fill=black!8},
      latente/.style  = {draw=black!75, fill=black!18, line width=0.6pt},
      texto/.style    = {font=\footnotesize, align=center},
      nombre/.style   = {font=\small, align=center},
      flecha/.style   = {->, draw=black!75, line width=0.7pt},
    ]
    \useasboundingbox (-1.7,-4.5) rectangle (12.2,4.6);

    %------------------------------------------------ vector de entrada
    \draw[vector] (0,-2.2) rectangle (0.5,2.2);
    \foreach \y in {-1.65,-1.1,...,1.75}
      \draw[celda] (0,\y) -- (0.5,\y);
    \node[nombre, anchor=south] at (0.25,2.6) {$\mathbf{x}\in\mathbb{R}^{n}$};
    \node[texto, anchor=north] at (0.25,-2.6) {entrada};

    %------------------------------------------------ encoder
    \draw[encoder] (1.4,-2.2) -- (4.3,-0.75) -- (4.3,0.75) -- (1.4,2.2) -- cycle;
    \node[nombre] at (2.85,0) {$E_{\phi}$};
    \node[texto, anchor=north] at (2.85,-2.6) {codificador};

    %------------------------------------------------ espacio latente
    \draw[latente] (5.2,-0.75) rectangle (5.7,0.75);
    \foreach \y in {-0.375,0,0.375}
      \draw[celda] (5.2,\y) -- (5.7,\y);
    \node[nombre, anchor=south] at (5.45,1.15) {$\mathbf{z}\in\mathbb{R}^{d}$};
    \node[texto, anchor=north] at (5.45,-1.15)
      {embedding\\[-1pt] $d\ll n$};

    %------------------------------------------------ decoder
    \draw[decoder] (6.6,-0.75) -- (9.5,-2.2) -- (9.5,2.2) -- (6.6,0.75) -- cycle;
    \node[nombre] at (8.05,0) {$D_{\theta}$};
    \node[texto, anchor=north] at (8.05,-2.6) {decodificador};

    %------------------------------------------------ vector de salida
    \draw[vector] (10.4,-2.2) rectangle (10.9,2.2);
    \foreach \y in {-1.65,-1.1,...,1.75}
      \draw[celda] (10.4,\y) -- (10.9,\y);
    \node[nombre, anchor=south] at (10.65,2.6) {$\hat{\mathbf{x}}\in\mathbb{R}^{n}$};
    \node[texto, anchor=north] at (10.65,-2.6) {reconstrucci\'on};

    %------------------------------------------------ flujo
    \draw[flecha] (0.6,0)  -- (1.3,0);
    \draw[flecha] (4.4,0)  -- (5.1,0);
    \draw[flecha] (5.8,0)  -- (6.5,0);
    \draw[flecha] (9.6,0)  -- (10.3,0);

    %------------------------------------------------ llaves de las etapas
    \draw[decorate, decoration={brace, amplitude=5pt}, draw=black!60]
      (1.4,2.75) -- (4.3,2.75);
    \node[texto, anchor=south] at (2.85,3.15)
      {compresi\'on\\[-1pt] $n\to d$};
    \draw[decorate, decoration={brace, amplitude=5pt}, draw=black!60]
      (6.6,2.75) -- (9.5,2.75);
    \node[texto, anchor=south] at (8.05,3.15)
      {expansi\'on\\[-1pt] $d\to n$};

    %------------------------------------------------ funcion de perdida
    \node[texto, anchor=north] at (5.45,-3.35)
      {$\mathcal{L}(\phi,\theta)=\lVert\mathbf{x}-\hat{\mathbf{x}}\rVert^{2}$};
  \end{tikzpicture}
  \caption{Esquema del autoencoder: el codificador comprime la entrada en un
    embedding de baja dimensión y el decodificador reconstruye la entrada a
    partir de él.}
  \label{fig:autoencoder-embedding}
\end{figure}
.
% Requiere en el preambulo del documento principal:
%   \usepackage{tikz}
%   \usetikzlibrary{decorations.pathreplacing,arrows.meta}
%   \usepackage{amsmath,amssymb}
%
% x     : vector de entrada de dimension n.
% z     : embedding (representacion latente) de dimension d << n.
% xhat  : reconstruccion de la entrada.
% E, D  : encoder y decoder, cada uno una MLP que reduce/expande la dimension.

\begin{figure}[htbp]
  \centering
  \begin{tikzpicture}[
      >={Latex[length=2.2mm]},
      vector/.style   = {draw=black!75, fill=white, line width=0.6pt},
      celda/.style    = {draw=black!35, line width=0.3pt},
      bloque/.style   = {draw=black!75, line width=0.6pt},
      encoder/.style  = {bloque, fill=black!8},
      decoder/.style  = {bloque, fill=black!8},
      latente/.style  = {draw=black!75, fill=black!18, line width=0.6pt},
      texto/.style    = {font=\footnotesize, align=center},
      nombre/.style   = {font=\small, align=center},
      flecha/.style   = {->, draw=black!75, line width=0.7pt},
    ]
    \useasboundingbox (-1.7,-4.5) rectangle (12.2,4.6);

    %------------------------------------------------ vector de entrada
    \draw[vector] (0,-2.2) rectangle (0.5,2.2);
    \foreach \y in {-1.65,-1.1,...,1.75}
      \draw[celda] (0,\y) -- (0.5,\y);
    \node[nombre, anchor=south] at (0.25,2.6) {$\mathbf{x}\in\mathbb{R}^{n}$};
    \node[texto, anchor=north] at (0.25,-2.6) {entrada};

    %------------------------------------------------ encoder
    \draw[encoder] (1.4,-2.2) -- (4.3,-0.75) -- (4.3,0.75) -- (1.4,2.2) -- cycle;
    \node[nombre] at (2.85,0) {$E_{\phi}$};
    \node[texto, anchor=north] at (2.85,-2.6) {codificador};

    %------------------------------------------------ espacio latente
    \draw[latente] (5.2,-0.75) rectangle (5.7,0.75);
    \foreach \y in {-0.375,0,0.375}
      \draw[celda] (5.2,\y) -- (5.7,\y);
    \node[nombre, anchor=south] at (5.45,1.15) {$\mathbf{z}\in\mathbb{R}^{d}$};
    \node[texto, anchor=north] at (5.45,-1.15)
      {embedding\\[-1pt] $d\ll n$};

    %------------------------------------------------ decoder
    \draw[decoder] (6.6,-0.75) -- (9.5,-2.2) -- (9.5,2.2) -- (6.6,0.75) -- cycle;
    \node[nombre] at (8.05,0) {$D_{\theta}$};
    \node[texto, anchor=north] at (8.05,-2.6) {decodificador};

    %------------------------------------------------ vector de salida
    \draw[vector] (10.4,-2.2) rectangle (10.9,2.2);
    \foreach \y in {-1.65,-1.1,...,1.75}
      \draw[celda] (10.4,\y) -- (10.9,\y);
    \node[nombre, anchor=south] at (10.65,2.6) {$\hat{\mathbf{x}}\in\mathbb{R}^{n}$};
    \node[texto, anchor=north] at (10.65,-2.6) {reconstrucci\'on};

    %------------------------------------------------ flujo
    \draw[flecha] (0.6,0)  -- (1.3,0);
    \draw[flecha] (4.4,0)  -- (5.1,0);
    \draw[flecha] (5.8,0)  -- (6.5,0);
    \draw[flecha] (9.6,0)  -- (10.3,0);

    %------------------------------------------------ llaves de las etapas
    \draw[decorate, decoration={brace, amplitude=5pt}, draw=black!60]
      (1.4,2.75) -- (4.3,2.75);
    \node[texto, anchor=south] at (2.85,3.15)
      {compresi\'on\\[-1pt] $n\to d$};
    \draw[decorate, decoration={brace, amplitude=5pt}, draw=black!60]
      (6.6,2.75) -- (9.5,2.75);
    \node[texto, anchor=south] at (8.05,3.15)
      {expansi\'on\\[-1pt] $d\to n$};

    %------------------------------------------------ funcion de perdida
    \node[texto, anchor=north] at (5.45,-3.35)
      {$\mathcal{L}(\phi,\theta)=\lVert\mathbf{x}-\hat{\mathbf{x}}\rVert^{2}$};
  \end{tikzpicture}
  \caption{Esquema del autoencoder: el codificador comprime la entrada en un
    embedding de baja dimensión y el decodificador reconstruye la entrada a
    partir de él.}
  \label{fig:autoencoder-embedding}
\end{figure}
